% -----------------------------------------------------------------------
% Related Work
% -----------------------------------------------------------------------

\paragraph{Inference serving frameworks.}
General-purpose inference servers provide multi-model hosting and hardware
acceleration but are designed for data-center workloads.
NVIDIA Triton Inference Server~\cite{nvidia2023triton} is highly capable,
supporting multiple backends and dynamic batching, yet its operational model
requires Docker or Kubernetes, Python deployment scripts, and a non-trivial
configuration surface that is impractical for edge or single-node
deployments.
TorchServe~\cite{pytorch2023torchserve} is tightly coupled to PyTorch,
requiring a Python environment for both model packaging and serving, and
provides no native support for ONNX-only runtimes.
BentoML~\cite{bentoml2023} offers a flexible Python-based abstraction for
packaging and deploying ML services, but its cloud-centric design and
Python runtime dependency make it a poor fit for lean, offline edge deployments.
Ollama~\cite{ollama2023} takes a philosophically different approach: a single
compiled binary that manages model downloads and serves inference over a local
REST API with zero external dependencies.
VisionServe adopts the same user-facing philosophy as Ollama --- one binary,
no account, no cloud, no Python runtime --- and applies it to the computer
vision domain, where the multi-session, multi-task, and execution-provider
complexity of CV models demands a purpose-built design.
None of the frameworks above provide a license-safe, binary-only serving
solution for the full range of modern CV tasks.

\paragraph{Computer vision model ecosystems.}
Ultralytics YOLO~\cite{ultralytics2023yolo} is among the most widely used
real-time object detectors, but it is released under the GNU Affero General
Public License v3.0 (AGPL-3.0)~\cite{agpllicense}.
AGPL is a strong copyleft: any software that incorporates or links against an
AGPL component must itself be distributed under AGPL, including commercial and
closed-source products.
This makes Ultralytics YOLO incompatible with permissive-license projects and
with any proprietary downstream deployment.
Torchvision~\cite{torchvision2023} ships BSD-3 licensed weights for standard
architectures but is a library, not a server; inference requires importing the
PyTorch Python stack.
HuggingFace Hub aggregates a large number of CV models with convenient
download APIs, but inference requires the Python \texttt{transformers} or
\texttt{diffusers} stack, and the host license of a Hub-hosted model is
determined by the uploader, not by the Hub itself --- a model's presence on
Hub says nothing about its license.
VisionServe addresses this gap by maintaining an explicit license allowlist
(Apache-2.0, MIT, and BSD families) validated at startup; models with any
other license are rejected outright, giving downstream adopters a strong
audit guarantee.

\paragraph{ONNX and cross-framework portability.}
The Open Neural Network Exchange format~\cite{onnx2019} defines a
framework-agnostic computation graph that enables models trained in PyTorch,
PaddlePaddle, TensorFlow, and other frameworks to be exported once and
deployed anywhere.
ONNX Runtime (ORT)~\cite{onnxruntime} is Microsoft's high-performance
inference engine for ONNX graphs.
ORT ships optional execution providers (EPs) for NVIDIA TensorRT and CUDA,
Apple CoreML, Intel OpenVINO, and DirectML for Windows GPUs, all behind a
common session API.
VisionServe builds on the \texttt{yalue/onnxruntime\_go} Go binding and
exposes ORT's EP selection through a manifest-level priority list, so a single
binary automatically dispatches to the best available accelerator at runtime.

\paragraph{Summary.}
Table~\ref{tab:related} contrasts the systems discussed above along the axes
most relevant to edge and embedded CV deployment.
VisionServe is unique in combining (i)~a single compiled binary with no
Python or container runtime, (ii)~strict permissive-license enforcement for
both the serving framework and every hosted model, (iii)~multi-EP hardware
fallback through ONNX Runtime, and (iv)~native support for prompted and
multi-session CV pipelines (segmentation, grounded detection) within a
unified result schema.

\begin{table}[h]
  \centering
  \caption{Comparison of inference serving systems.}
  \label{tab:related}
  \small
  \begin{tabular}{lccccc}
    \toprule
    \textbf{System} & \textbf{No Python} & \textbf{Single binary} &
    \textbf{License audit} & \textbf{Multi-EP} & \textbf{CV-native} \\
    \midrule
    Triton~\cite{nvidia2023triton}        & \texttimes & \texttimes & \texttimes & \checkmark & \texttimes \\
    TorchServe~\cite{pytorch2023torchserve} & \texttimes & \texttimes & \texttimes & \texttimes & \texttimes \\
    BentoML~\cite{bentoml2023}            & \texttimes & \texttimes & \texttimes & \texttimes & \texttimes \\
    Ollama~\cite{ollama2023}              & \checkmark & \checkmark & \texttimes & \texttimes & \texttimes \\
    \midrule
    \textbf{VisionServe} (ours)           & \checkmark & \checkmark & \checkmark & \checkmark & \checkmark \\
    \bottomrule
  \end{tabular}
\end{table}
